\documentclass[11pt,a4paper]{article}

% Pakete
\usepackage[utf8]{inputenc}
\usepackage[T1]{fontenc}
\usepackage[ngerman]{babel}
\usepackage[left=18mm,right=18mm,top=18mm,bottom=18mm]{geometry}
\usepackage{helvet}
\renewcommand{\familydefault}{\sfdefault}
\usepackage{amsmath,amssymb}
\usepackage{tikz}
\usetikzlibrary{arrows.meta,positioning,decorations.pathmorphing}
\usepackage{tcolorbox}
\tcbuselibrary{skins,breakable}
\usepackage{enumitem}
\usepackage{tabularx}
\usepackage{siunitx}
\sisetup{locale=DE, per-mode=fraction}
\usepackage{qrcode}

% Fachfarbe Physik
\definecolor{physik}{HTML}{2c5aa0}
\definecolor{physiklight}{HTML}{e3f2fd}
\definecolor{erfolg}{HTML}{2a7a4b}
\definecolor{erfolglight}{HTML}{e8f5e9}

% Box-Stile
\tcbset{
    infobox/.style={
        colback=physiklight,
        colframe=physik,
        boxrule=1pt,
        arc=2pt,
        left=6pt, right=6pt, top=5pt, bottom=5pt
    },
    schrittbox/.style={
        colback=white,
        colframe=physik,
        fonttitle=\bfseries,
        title=#1,
        boxrule=1.5pt,
        arc=2pt,
        left=6pt, right=6pt, top=5pt, bottom=5pt
    },
    ergebnisbox/.style={
        colback=erfolglight,
        colframe=erfolg,
        boxrule=2pt,
        arc=2pt,
        left=8pt, right=8pt, top=6pt, bottom=6pt
    }
}

% Kopfzeile
\usepackage{fancyhdr}
\pagestyle{fancy}
\fancyhf{}
\lhead{\small Physik 12 GK}
\chead{\small Lernblatt}
\rhead{\small 14.01.2026}
\renewcommand{\headrulewidth}{0.4pt}

% Keine Einrückung
\setlength{\parindent}{0pt}
\setlength{\parskip}{5pt}

\begin{document}

% ═══════════════════════════════════════════════════════════════════════════
% TITEL
% ═══════════════════════════════════════════════════════════════════════════

\begin{minipage}[t]{0.72\textwidth}
\vspace{0pt}
{\Large\textbf{Photoeffekt -- h-Bestimmung}}\\[0.2em]
{\normalsize Fortsetzung: Messung \& Auswertung (ES)}
\end{minipage}
\hfill
\begin{minipage}[t]{0.25\textwidth}
\vspace{0pt}
\raggedleft
\qrcode[height=1.8cm]{https://devbox42.github.io/sciencesim/physik/kl12-GK/01-photoeffekt/02-DS-gegenfeld/SIM-02-gegenfeld.html}\\[2pt]
{\scriptsize Simulation}
\end{minipage}

\vspace{0.4em}
\hrule height 1.5pt
\vspace{0.6em}

% ═══════════════════════════════════════════════════════════════════════════
% WIEDERHOLUNG (kompakt)
% ═══════════════════════════════════════════════════════════════════════════

\begin{tcolorbox}[infobox]
\textbf{Letzte Stunde:} Gegenfeldmethode -- Ein elektrisches Feld bremst die Photoelektronen ab.

\vspace{0.3em}
\begin{minipage}[t]{0.55\textwidth}
\textbf{Grenzspannung} $U_G$ stoppt die schnellsten Elektronen:
\begin{center}
$E_{\text{kin,max}} = e \cdot U_G$
\end{center}
\textbf{Praktisch:} Bei $U_G$ in Volt ist $E_{\text{kin}}$ numerisch gleich in \textbf{eV}.
\end{minipage}
\hfill
\begin{minipage}[t]{0.40\textwidth}
\centering
\begin{tikzpicture}[scale=0.55]
    % Fotozelle kompakt
    \draw[thick] (0,0) rectangle (4,2.8);
    \draw[very thick, physik] (0.6,0.3) -- (0.6,2.5);
    \draw[very thick] (3.4,0.3) -- (3.4,2.5);
    \node[left, font=\scriptsize] at (0.5,1.4) {K};
    \node[right, font=\scriptsize] at (3.5,1.4) {A};
    % Licht
    \draw[->, thick, yellow!80!orange, decorate, decoration={snake, amplitude=1.5pt, segment length=6pt}] (-0.8,1.8) -- (0.5,1.8);
    % Elektronen
    \draw[->, thick, blue] (0.8,1.4) -- (1.6,1.4);
    \node[blue, font=\scriptsize] at (1.2,1.8) {$e^-$};
    % Gegenfeld
    \draw[<-, thick, red] (2,0.8) -- (3.2,0.8);
    % Spannung
    \draw (0.6,-0.2) -- (0.6,-0.6) -- (1.6,-0.6);
    \draw (3.4,-0.2) -- (3.4,-0.6) -- (2.4,-0.6);
    \draw (1.6,-0.8) rectangle (2.4,-0.4);
    \node[font=\scriptsize] at (2,-0.6) {$U_G$};
\end{tikzpicture}
\end{minipage}
\end{tcolorbox}

% ═══════════════════════════════════════════════════════════════════════════
% SCHRITT 1: MESSUNG
% ═══════════════════════════════════════════════════════════════════════════

\begin{tcolorbox}[schrittbox={Schritt 1: Messreihe aufnehmen (15 min)}]
\begin{enumerate}[nosep, leftmargin=*]
\item Öffne die \textbf{Simulation} (QR-Code oben rechts)
\item Wähle \textbf{Natrium} als Material
\item Für jede Wellenlänge in der Tabelle:
    \begin{itemize}[nosep]
    \item Wellenlänge einstellen (Wert direkt eintippen!)
    \item Batteriespannung erhöhen, bis \textbf{Strom = 0}
    \item $U_G$ ablesen und in \textbf{AB-02, Aufgabe 1} eintragen
    \end{itemize}
\item Bei manchen Farben: \textbf{kein Photoeffekt} $\rightarrow$ ,,---'' eintragen
\end{enumerate}

\vspace{0.3em}
\textit{Tipp: Wellenlänge und Spannung lassen sich feiner einstellen, wenn du den Wert direkt eintippst (Klick aufs Zahlenfeld).}
\end{tcolorbox}

% ═══════════════════════════════════════════════════════════════════════════
% SCHRITT 2: DIAGRAMM
% ═══════════════════════════════════════════════════════════════════════════

\begin{tcolorbox}[schrittbox={Schritt 2: Einsteinsche Gerade zeichnen (15 min)}]
\begin{enumerate}[nosep, leftmargin=*]
\item Trage deine Messwerte in das Diagramm ein (\textbf{AB-02, Aufgabe 2}):
    \begin{itemize}[nosep]
    \item x-Achse: Frequenz $f$ (in $10^{14}$ Hz), Skala: 0--8
    \item y-Achse: $E_{\text{kin}}$ (in eV), Skala: $-4$ bis $+2$
    \end{itemize}
\item Zeichne eine \textbf{Ausgleichsgerade} durch die Messpunkte
\item \textbf{Extrapoliere} die Gerade nach links bis zur y-Achse!
\end{enumerate}

\vspace{0.3em}
\textbf{Beobachtung:} Die Punkte liegen auf einer \underline{\hspace{2.5cm}} $\rightarrow$ linearer Zusammenhang!
\end{tcolorbox}

% ═══════════════════════════════════════════════════════════════════════════
% SCHRITT 3: AUSWERTUNG
% ═══════════════════════════════════════════════════════════════════════════

\begin{tcolorbox}[schrittbox={Schritt 3: Auswertung (15 min)}]
\textbf{a) Steigung bestimmen} (\textbf{AB-02, Aufgabe 3a})

Wähle zwei Punkte auf der Geraden und berechne:
\begin{center}
$m = \dfrac{\Delta E}{\Delta f} = \dfrac{E_2 - E_1}{f_2 - f_1}$
\end{center}

\textbf{Die Steigung ist das Plancksche Wirkungsquantum $h$!} (Literaturwert: $h = \SI{4,14e-15}{eV.s}$)

\vspace{0.4em}
\textbf{b) Achsenabschnitte ablesen} (\textbf{AB-02, Aufgabe 3b})

\begin{tabular}{@{}p{7cm}l@{}}
x-Achsenabschnitt $f_G$ (wo $E_{\text{kin}} = 0$): & = \textbf{Grenzfrequenz} \\
y-Achsenabschnitt (bei $f = 0$): & = $-W_A$ (negative \textbf{Austrittsarbeit})
\end{tabular}

\vspace{0.4em}
\textbf{c) Austrittsarbeit berechnen} (\textbf{AB-02, Aufgabe 3c})

$W_A = h \cdot f_G$ \quad oder \quad $W_A = |$y-Achsenabschnitt$|$ \quad (Literaturwert Natrium: $W_A = \SI{2,3}{eV}$)
\end{tcolorbox}

\vspace{0.5em}

% ═══════════════════════════════════════════════════════════════════════════
% ERGEBNIS
% ═══════════════════════════════════════════════════════════════════════════

\begin{tcolorbox}[ergebnisbox]
\textbf{Ergebnis: Die Photogleichung}

\vspace{0.3em}
\begin{center}
{\Large $\boxed{E_{\text{kin,max}} = h \cdot f - W_A}$}
\end{center}

\vspace{0.3em}
\begin{minipage}[t]{0.48\textwidth}
\textbf{Die Gerade zeigt:}
\begin{itemize}[nosep]
\item Steigung = $h$ (Naturkonstante!)
\item x-Achsenabschnitt = $f_G$
\item y-Achsenabschnitt = $-W_A$
\end{itemize}
\end{minipage}
\hfill
\begin{minipage}[t]{0.48\textwidth}
\textbf{Physikalische Bedeutung:}
\begin{itemize}[nosep]
\item $h \cdot f$ = Photonenenergie
\item $W_A$ = Energie zum Herauslösen
\item $f < f_G$: kein Photoeffekt
\end{itemize}
\end{minipage}
\end{tcolorbox}

\vspace{0.5em}

% ═══════════════════════════════════════════════════════════════════════════
% REFLEXION (falls Zeit)
% ═══════════════════════════════════════════════════════════════════════════

\textbf{Falls Zeit: Reflexion} $\rightarrow$ \textbf{AB-02, Aufgabe 4}

\begin{itemize}[nosep]
\item Warum konnte $h$ aus dem Experiment bestimmt werden?
\item Welche Fehlerquellen gibt es bei deiner h-Bestimmung?
\end{itemize}

\vspace{0.3em}
\hrule height 0.5pt
\vspace{0.3em}

{\small\textbf{Zusammenfassung:} Du hast aus eigenen Messdaten das \textbf{Plancksche Wirkungsquantum} $h$ bestimmt -- eine fundamentale Naturkonstante der Quantenphysik! Einstein erhielt dafür 1921 den Nobelpreis.}

\end{document}
